\documentclass[11pt,a4paper]{article}
\usepackage[utf8]{inputenc}
\usepackage[T5]{fontenc}
\usepackage[vietnamese]{babel}
\usepackage{amsmath,amssymb,amsfonts}
\usepackage{booktabs}
\usepackage{array}
\usepackage[margin=2.4cm]{geometry}
\usepackage{enumitem}
\usepackage{xcolor}

% ---- Ký hiệu chuẩn hóa (macro) ------------------------------------------
\newcommand{\Real}{\mathbb{R}}
\newcommand{\Integer}{\mathbb{Z}}
\newcommand{\Iset}{\mathcal{I}}      % tập cọc
\newcommand{\Kset}{\mathcal{K}}      % tập tổ hợp tải
\newcommand{\Tset}{\mathcal{T}}      % tập kiểu lưới
\newcommand{\Dset}{\mathcal{D}}      % tập đường kính
\newcommand{\Hset}{\mathcal{H}}      % ràng buộc cứng
\newcommand{\Sset}{\mathcal{S}}      % ràng buộc mềm
\newcommand{\Feas}{\Omega}           % miền khả thi
\newcommand{\pos}[1]{\left[#1\right]^{+}}  % phần dương
\newcommand{\Pall}{\overline{P}}     % [Po]
\newcommand{\Tall}{\overline{T}}     % [Ct]
\newcommand{\Mall}{\overline{M}}     % [M]
\newcommand{\Hall}{\overline{H}}     % [H]
\newcommand{\Pmax}{P_{\max}}
\newcommand{\Pmin}{P_{\min}}
\newcommand{\Mxx}{M^{x}_{\max}}
\newcommand{\Myy}{M^{y}_{\max}}
\newcommand{\Hmx}{H_{\max}}
\newcommand{\xvec}{\mathbf{x}}
\newcommand{\Gen}{\mathbf{G}}        % toán tử sinh lưới
\newcommand{\Mcoc}{\boldsymbol{\Psi}}% toán tử MCOC

\title{\textbf{Mô hình tối ưu bố trí cọc móng cầu}\\[2pt]
\large Quy hoạch đa mục tiêu có ràng buộc cứng -- mềm}
\author{OptApp --- Mô tả toán học}
\date{}

\begin{document}
\maketitle

\section{Bảng ký hiệu (chuẩn hóa)}

\subsection*{Tập hợp và chỉ số}
\begin{tabular}{@{}ll@{}}
\toprule
Ký hiệu & Ý nghĩa \\
\midrule
$i \in \Iset=\{1,\dots,n\}$ & chỉ số cọc \\
$k \in \Kset=\{1,\dots,K\}$ & chỉ số tổ hợp tải trọng \\
$t \in \Tset=\{\mathrm{A},\mathrm{B}\}$ & kiểu lưới: A trực giao, B hoa mai (so le) \\
$d \in \Dset=\{D_1,\dots,D_m\}$ & tập đường kính cọc (mô hình mở rộng) \\
$r \in \Hset$, \; $q \in \Sset$ & chỉ số ràng buộc cứng / mềm \\
\bottomrule
\end{tabular}

\subsection*{Tham số đầu vào (cho trước)}
\begin{tabular}{@{}lll@{}}
\toprule
Ký hiệu & Ý nghĩa & Đơn vị \\
\midrule
$D$ & đường kính cọc & m \\
$L_x,\,L_y$ & kích thước bệ (cap) theo $x,y$ & m \\
$c$ & khoảng cách an toàn tim cọc tới mép bệ (SAFE\_D) & m \\
$\alpha_\ell=3,\ \alpha_u=6$ & hệ số khoảng cách tim cọc tối thiểu / tối đa & -- \\
$\Pall$ & sức chịu nén cho phép $[P_o]$ (hoặc $R_{c,d}$ theo TCVN) & T \\
$\Tall$ & sức chịu nhổ cho phép $[C_t]$ & T \\
$\Mall$ & mômen uốn cho phép $[M]$ & T$\cdot$m \\
$\Hall$ & lực ngang cho phép $[H]$ & T \\
$N_k,H^x_k,H^y_k,M^x_k,M^y_k,M^z_k$ & thành phần tổ hợp tải $k$ & T, T$\cdot$m \\
$\gamma_0,\gamma_n,\gamma_k$ & hệ số an toàn TCVN 10304:2014 (Đ.7.1.11) & -- \\
\bottomrule
\end{tabular}

\subsection*{Biến quyết định}
\begin{tabular}{@{}lll@{}}
\toprule
Ký hiệu & Miền & Ý nghĩa \\
\midrule
$t$ & $\Tset$ & kiểu lưới \\
$n_x$ & $\{1,\dots,n_x^{\max}\}\subset\Integer$ & số cột cọc \\
$n_y$ & $\{1,\dots,n_y^{\max}\}\subset\Integer$ & số hàng cọc \\
$s_x$ & $[\alpha_\ell D,\ \alpha_u D]\subset\Real$ & bước lưới theo $x$ \\
$s_y$ & $[\alpha_\ell D,\ \alpha_u D]\subset\Real$ & bước lưới theo $y$ \\
\bottomrule
\end{tabular}

\medskip
\noindent Véc-tơ quyết định: $\xvec=(t,n_x,n_y,s_x,s_y)$ (mô hình mở rộng: $\xvec=(d,t,n_x,n_y,s_x,s_y)$).

\subsection*{Biến trạng thái (dẫn xuất từ $\xvec$)}
\begin{tabular}{@{}ll@{}}
\toprule
Ký hiệu & Ý nghĩa \\
\midrule
$\mathbf{p}_i=(x_i,y_i)$ & tọa độ tim cọc $i$ \\
$\mathcal{C}=\{\mathbf{p}_i\}_{i\in\Iset}$ & tập tọa độ; $n=|\mathcal{C}|$ là số cọc \\
$s^{\mathrm{diag}}=\sqrt{(s_x/2)^2+s_y^2}$ & khoảng cách đường chéo (chỉ kiểu B) \\
$P_i^{k}$ & phản lực dọc trục cọc $i$ ở tổ hợp $k$ (dương = nén) \\
$\Pmax,\Pmin$ & lực nén lớn nhất / lực kéo lớn nhất trên mọi $(i,k)$ \\
$\Mxx,\Myy$ & mômen đầu cọc lớn nhất quanh $x,y$ \\
$\Hmx$ & lực ngang tổng hợp lớn nhất trên một cọc \\
$\Phi(\mathcal{C})$ & diện tích bao (footprint) của nhóm cọc \\
\bottomrule
\end{tabular}

\section{Toán tử mô hình}

\paragraph{Sinh tọa độ.} Toán tử $\Gen:\xvec\mapsto\mathcal{C}$ đặt lưới đối xứng quanh tâm bệ:
\begin{equation}
x_\iota=\Big(\iota-\tfrac{n_x-1}{2}\Big)s_x,\quad
y_\jmath=\Big(\jmath-\tfrac{n_y-1}{2}\Big)s_y,\qquad
\iota=0,\dots,n_x-1,\ \jmath=0,\dots,n_y-1,
\end{equation}
với kiểu B các hàng lẻ giảm một cọc và tự lệch $s_x/2$ (bố trí hoa mai).

\paragraph{Đánh giá nội lực (oracle MCOC).} Toán tử chính xác $\Mcoc$ trả phản lực cọc:
\begin{equation}
\big\{P_i^{k}\big\}_{i\in\Iset,\,k\in\Kset}=\Mcoc\!\big(\mathcal{C};\,\{N_k,H^x_k,H^y_k,M^x_k,M^y_k,M^z_k\}_{k\in\Kset}\big),
\end{equation}
\begin{equation}
\Pmax=\max_{i,k}P_i^{k},\qquad \Pmin=\min_{i,k}P_i^{k}.
\end{equation}
$\Mcoc$ là \emph{hộp đen} (không khả vi, không lồi) $\Rightarrow$ bài toán phải giải bằng metaheuristic.

\section{Hàm mục tiêu}

Bài toán \textbf{đa mục tiêu}, cùng cực tiểu hóa hai mục tiêu:
\begin{equation}
\min_{\xvec\in\Feas}\ \mathbf{F}(\xvec)=\big(f_1(\xvec),\,f_2(\xvec)\big),
\end{equation}
\begin{equation}
f_1(\xvec)=n \quad\text{(số cọc --- chi phí vật liệu \& thi công)},
\end{equation}
\begin{equation}
f_2(\xvec)=
\begin{cases}
\Phi(\mathcal{C}) & \text{chế độ ``bệ gọn'' (tiết kiệm bê tông bệ)},\\[3pt]
\Pmax & \text{chế độ ``an toàn'' (tăng dự trữ chịu lực)}.
\end{cases}
\end{equation}
Nghiệm là \textbf{tập Pareto} chứ không phải một điểm; mục tiêu phụ $f_2$ do người dùng chọn.

\paragraph{Mục tiêu mở rộng.} Khi đường kính là biến, hàm mục tiêu thay bằng \emph{chi phí vật liệu} (xấp xỉ thể tích bê tông cọc trên một mét dài):
\begin{equation}
\min_{d\in\Dset}\ \mathcal{V}(\xvec)=n\cdot\frac{\pi d^{2}}{4},
\end{equation}
đồng hạng thì ưu tiên ít cọc hơn ($\min n$).

\section{Ràng buộc cứng (hard constraints)}
Ràng buộc cứng \emph{bắt buộc} thỏa; vi phạm $\Rightarrow$ phương án \emph{bất khả thi}. Mỗi ràng buộc viết dạng chuẩn $g_r(\xvec)\le 0$.

\begin{equation}
\begin{aligned}
\text{(H1) Tồn tại cọc:}\quad & g_1 = 1-n \le 0,\\
\text{(H2) Khoảng cách tối thiểu (TCVN, BẮT BUỘC):}\quad & g_2 = \alpha_\ell D - s^{\star} \le 0,\\
\text{(H3) Mép bệ theo $x$:}\quad & g_3 = \max_i|x_i|+c-\tfrac{L_x}{2}\le 0,\\
\text{(H4) Mép bệ theo $y$:}\quad & g_4 = \max_i|y_i|+c-\tfrac{L_y}{2}\le 0,\\
\text{(H5) Sức chịu nén:}\quad & g_5 = \Pmax-\Pall \le 0,\\
\text{(H6) Sức chịu nhổ ($\Tall>0$):}\quad & g_6 = -\Tall-\Pmin \le 0,\\
\text{(H7) Mômen uốn ($\Mall>0$):}\quad & g_7 = \max(\Mxx,\Myy)-\Mall \le 0,\\
\text{(H8) Lực ngang ($\Hall>0$):}\quad & g_8 = \Hmx-\Hall \le 0,\\
\text{(H9) Tương tác $P$--$M$:}\quad & g_9 = \dfrac{\Pmax}{\Pall}+\dfrac{\max(\Mxx,\Myy)}{\Mall}-1 \le 0.
\end{aligned}
\end{equation}
trong đó $s^\star=\min\{s_x,s_y\}$ với kiểu A, và $s^\star=\min\{s_x,\,s^{\mathrm{diag}}\}$ với kiểu B (cặp gần nhất là đường chéo). H8, H9 được kích hoạt ở luồng mở rộng (R7/R8).

\paragraph{Miền khả thi.}
\begin{equation}
\Feas=\big\{\xvec:\ g_r(\xvec)\le 0,\ \forall r\in\Hset\big\},\qquad
\Hset=\{1,\dots,9\}.
\end{equation}

\section{Ràng buộc mềm (soft constraints)}
Ràng buộc mềm \emph{không loại} phương án mà bị \emph{phạt} (preference); biểu diễn dạng $\tilde{g}_q(\xvec)\le 0$.

\begin{equation}
\begin{aligned}
\text{(S1) Khoảng cách tối đa $\le 6D$:}\quad & \tilde{g}_1 = s^{\dagger}-\alpha_u D \le 0,\\
\text{(S2) Tính đều/đối xứng lưới:}\quad & \text{áp đặt cấu trúc bởi }\Gen\ (\text{không phạt thêm}),
\end{aligned}
\end{equation}
với $s^\dagger=\max\{s_x,s_y\}$ (kiểu A) hoặc $\max\{s_x,s^{\mathrm{diag}}\}$ (kiểu B). S1 là \emph{quy ước thực hành}: chỉ kích hoạt khi cờ \texttt{ENFORCE\_SPACING\_MAX} bật, ngược lại chỉ cảnh báo.

\section{Vô hướng hóa vi phạm và quan hệ trội}

\paragraph{Vi phạm cứng (chuẩn hóa không thứ nguyên).}
\begin{equation}
\Theta(\xvec)=\sum_{r\in\Hset}\frac{\pos{g_r(\xvec)}}{\rho_r},\qquad
\pos{a}=\max(0,a),
\end{equation}
với $\rho_r$ là đại lượng giới hạn tương ứng ($\rho_5=\Pall$, $\rho_3=L_x/2,\dots$) đưa mọi thành phần về cùng thang $\mathcal{O}(1)$. Phương án \textbf{khả thi} $\iff \Theta(\xvec)=0$.

\paragraph{Phạt mềm (gộp vào mục tiêu phụ).}
\begin{equation}
\Pi(\xvec)=\sum_{q\in\Sset} w_q\,\frac{\pos{\tilde{g}_q(\xvec)}}{\tilde{\rho}_q},\qquad
f_2^{\text{aug}}(\xvec)=f_2(\xvec)\big(1+\Pi(\xvec)\big),
\end{equation}
với $w_q\ge 0$ trọng số ưu tiên. Khi mọi ràng buộc mềm thỏa thì $\Pi=0$ và $f_2^{\text{aug}}=f_2$.

\paragraph{Quan hệ trội có ràng buộc (constrained-domination).}
$\xvec^{(a)}\prec\xvec^{(b)}$ khi và chỉ khi
\begin{equation}
\begin{cases}
\Theta_a=0\ \wedge\ \Theta_b>0, & \text{(khả thi thắng bất khả thi)}\\[3pt]
0<\Theta_a<\Theta_b, & \text{(ít vi phạm cứng hơn thắng)}\\[3pt]
\Theta_a=\Theta_b=0\ \wedge\ \mathbf{F}^{\text{aug}}(\xvec^{(a)})\preceq_{\mathrm{Pareto}}\mathbf{F}^{\text{aug}}(\xvec^{(b)}), & \text{(cả hai khả thi $\to$ Pareto)}
\end{cases}
\end{equation}
với $\mathbf{F}^{\text{aug}}=(f_1,f_2^{\text{aug}})$. Tập nghiệm là
$\mathcal{P}=\{\xvec\in\Feas:\nexists\,\xvec'\in\Feas,\ \xvec'\prec\xvec\}$.

\section{Phát biểu gọn của mô hình}

\begin{equation}
\boxed{
\begin{aligned}
&\min_{\xvec}\quad \mathbf{F}^{\text{aug}}(\xvec)=\Big(\,n,\ \ f_2(\xvec)\,[1+\Pi(\xvec)]\,\Big)\\[4pt]
\text{(cứng)}\quad &\text{s.t.}\quad g_r(\xvec)\le 0,\quad r\in\Hset=\{1,\dots,9\}\\[2pt]
\text{(mềm)}\quad &\phantom{\text{s.t.}}\quad \tilde{g}_q(\xvec)\le 0\ \text{(phạt, không bắt buộc)},\quad q\in\Sset\\[2pt]
&\phantom{\text{s.t.}}\quad t\in\Tset,\ n_x,n_y\in\Integer_{+},\ s_x,s_y\in[\alpha_\ell D,\alpha_u D]\\[2pt]
&\phantom{\text{s.t.}}\quad \mathcal{C}=\Gen(\xvec),\ \{P_i^k\}=\Mcoc(\mathcal{C};\{\ell_k\})
\end{aligned}}
\end{equation}

\paragraph{Phương pháp giải.} Vì $\Mcoc$ là hộp đen không khả vi, bài toán đa mục tiêu, biến hỗn hợp, miền khả thi rất hẹp $\Rightarrow$ dùng \textbf{NSGA-II} (sắp xếp không-bị-trội nhanh + khoảng cách chen chúc + SBX/đột biến đa thức + chọn lọc $\mu{+}\lambda$), kèm \emph{gieo hạt tất định} (phủ vùng khả thi), \emph{cache} theo spec lưới và \emph{trần ngân sách} số lần gọi $\Mcoc$.

\end{document}
